\section{2026 年 11 月核验增补}

本节按论文首次提交时间排列。每项方法附有可核验的外部入选依据。

% METHOD: 2607.08769
\paragraph{PanoLOG：将大范围全景重建分为全局粗建与局部细化}
2026.07~\cite{2607_08769}。PanoLOG 先全局粗训练，再按几何结构划分空间、按相机对各块的梯度贡献分配训练图像，以缓解全景相机使过多区域同时可见的问题。模型同时使用显式天空球和全景单目深度监督，帮助约束无界背景和观测困难区域。

作者以 Pano360 进行主体量化和渲染对照，并在 360Roam、Ricoh360 追加可视化。论文假设场景静止，未抑制的车辆和行人可能留下重影，初始全局训练的大规模内存要求也仍存在。作者单位明确列出 Insta360 Research。

\textbf{入选依据} 论文署名单位包括 Insta360 Research（\href{https://arxiv.org/html/2607.08769v1}{论文原文}）。

% METHOD: 2611.22805
\paragraph{DOOM-GS：以时空分块和流式泼溅训练城市级动态高斯}
2026.11~\cite{2611_22805}。动态场景离开显存训练时，一个视角的可见原语也可能塞不进显存。DOOM-GS 按空间与时间组织数据块，利用相邻相机可见集合重叠来排程，并以流式泼溅和梯度估计在有界显存下更新模型。

论文在加入程序生成车辆和行人的 DynamicMatrixCity 比较单一动态模型与逐帧静态基线，并使用留出相机质量、静态区闪烁和训练时间衡量。结果针对该模拟城市场景和相机划分，真实交通动态的效果需另行验证。

\textbf{入选依据} 论文署名单位包括 Baidu（\href{https://arxiv.org/abs/2611.22805}{论文原文}）。
